\chapter{Preparacion de competencia}
\label{cap:preparacion-competencia}

\begin{procedimiento}{Objetivo del capitulo}{Operador del torneo}
Guiar la verificacion previa y la configuracion inicial de categorias antes de operar grupos, batallas o fases posteriores.
\end{procedimiento}

La preparacion de competencia inicia cuando registros, asistencia y equipos ya fueron revisados segun \cref{cap:gestion-organizacion}. Este capitulo cubre el estado inicial de una categoria, el formato sugerido, el formato manual y la distribucion de grupos. No documenta todavia seleccion de clasificados, resultados de batallas, repechaje ni podio.

\section{Verificaciones antes de inicializar}

\begin{tabularx}{\linewidth}{>{\bfseries}p{4.2cm}X}
\toprule
Revision & Criterio operativo\\
\midrule
Participantes & Los registros corresponden a la edicion activa y no contienen duplicados evidentes.\\
Asistencia & Los equipos que asistiran al evento estan confirmados o revisados por la organizacion.\\
Equipos & Cada robot tiene institucion, categoria e integrantes esperados.\\
Division & La categoria a preparar es Colegios o Universidades, sin mezclar registros.\\
Permisos & El operador inicio sesion como usuario autorizado y entiende el impacto de regenerar grupos.\\
Progreso existente & Antes de cambiar formato o volver a mezclar, se revisa si ya hay resultados o fases creadas.\\
\bottomrule
\end{tabularx}

\begin{advertencia}
No inicialice una categoria si los conteos no coinciden con la lista revisada por la organizacion. Primero corrija asistencia o equipos desde las superficies correspondientes y vuelva al panel cuando los datos esten listos.
\end{advertencia}

\section{Entrar a una categoria preparada}

\figuraManual{capturas/finales/MU-015-configuracion-division.png}{Resumen de una categoria con navegacion de fases, estado de competencia, modo activo y accion rapida para volver a mezclar grupos.}{fig:bloque2-mu-015}

\begin{procedimiento}{Revisar el estado de la categoria}{Operador del torneo}
Use este procedimiento despues de abrir \botoninterfaz{Abrir Colegios} o \botoninterfaz{Abrir Universidades} desde el panel.
\end{procedimiento}

\begin{precondiciones}
\begin{itemize}
    \item La categoria ya tiene tablero base generado.
    \item El operador esta en la pantalla de resumen de la categoria.
    \item La revision de asistencia y equipos fue realizada antes de esta operacion.
\end{itemize}
\end{precondiciones}

\paso{1}{Confirme el titulo de la categoria, por ejemplo \campointerfaz{Colegios}.}
\paso{2}{Revise la navegacion superior de fases. En esta etapa use \botoninterfaz{Resumen}; las fases posteriores se documentan mas adelante.}
\paso{3}{Verifique \campointerfaz{Edicion}, \campointerfaz{Estado} y \campointerfaz{Modo activo}.}
\paso{4}{Revise la semilla mostrada en \campointerfaz{Estado}; sirve para reconocer la mezcla actual de grupos.}
\paso{5}{Revise el panel \campointerfaz{Recomendacion del sistema} sin crear fases posteriores desde este capitulo.}
\paso{6}{Si debe reconstruir la distribucion inicial, use \botoninterfaz{Volver a mezclar grupos} solo despues de validar que no hay progreso que conservar.}

\begin{resultadoesperado}
El operador identifica si la categoria esta no inicializada, preparada o con progreso visible antes de realizar cualquier cambio.
\end{resultadoesperado}

\begin{fallas}
Si la categoria no abre, vuelva al panel y confirme que el tablero fue generado. Si el estado o el modo no coinciden con lo esperado, revise los conteos y no avance a fases posteriores.
\end{fallas}

\section{Estados previos de competencia}

\begin{tabularx}{\linewidth}{>{\bfseries}p{4cm}X}
\toprule
Estado operativo & Como actuar\\
\midrule
No inicializada & En el panel de organizacion aparece la opcion \botoninterfaz{Generar torneo}. Revise asistencia y equipos antes de crear el tablero.\\
Preparada & La categoria abre en resumen, muestra grupos y modo activo. Puede revisarse formato o distribucion antes de registrar resultados.\\
Con progreso existente & La navegacion muestra fases con actividad o resultados. No use acciones de mezcla o formato sin revisar el avance.\\
Recalculada o regenerada & La interfaz informa que aseguro o mezclo la categoria. Verifique grupos, participantes y modo antes de continuar.\\
\bottomrule
\end{tabularx}

\section{Aplicar modo sugerido}

\figuraManual{capturas/finales/MU-016-formato-manual.png}{Zona de configuracion inicial con modo sugerido y formulario editable para preparar la primera fase.}{fig:bloque2-mu-016}

\begin{procedimiento}{Usar el formato sugerido}{Operador del torneo}
Use esta opcion cuando la organizacion acepta la recomendacion inicial de la plataforma.
\end{procedimiento}

\begin{precondiciones}
\begin{itemize}
    \item El operador reviso la cantidad de equipos de la categoria.
    \item No existe progreso que deba preservarse con una distribucion diferente.
    \item La recomendacion coincide con las reglas operativas del evento.
\end{itemize}
\end{precondiciones}

\paso{1}{En la categoria, baje hasta \campointerfaz{Configuracion inicial sugerida}.}
\paso{2}{Presione \botoninterfaz{Ver detalle} si la seccion esta cerrada.}
\paso{3}{Revise cantidad de grupos, clasificados por grupo y si la configuracion incluye repechaje o tercer lugar.}
\paso{4}{Presione \botoninterfaz{Usar modo sugerido}.}
\paso{5}{Espere el mensaje de confirmacion y vuelva a revisar \campointerfaz{Modo activo}.}

\begin{resultadoesperado}
La plataforma muestra un mensaje como \campointerfaz{Se aplico el modo sugerido automatico sin borrar progreso existente.} y actualiza el resumen de la categoria.
\end{resultadoesperado}

\begin{fallas}
Si aparece un error de configuracion, revise cantidad de grupos y clasificados permitidos. No fuerce un formato manual hasta confirmar que la lista de equipos es correcta.
\end{fallas}

\section{Aplicar formato manual}

\begin{procedimiento}{Preparar primera fase con formato manual}{Operador del torneo}
Use esta opcion cuando la organizacion necesita ajustar grupos, clasificados, purgatorio o tercer lugar antes de iniciar.
\end{procedimiento}

\paso{1}{Abra \campointerfaz{Configuracion inicial editable} con \botoninterfaz{Abrir ajustes}.}
\paso{2}{Defina \campointerfaz{Numero de grupos}.}
\paso{3}{Defina \campointerfaz{Clasificados por grupo}.}
\paso{4}{Active o desactive \campointerfaz{Activar Purgatorio 1}, \campointerfaz{Activar Purgatorio 2} y \campointerfaz{Jugar tercer lugar} segun la decision de la organizacion.}
\paso{5}{Presione \botoninterfaz{Aplicar formato manual}.}
\paso{6}{Revise el mensaje de resultado y confirme que \campointerfaz{Modo activo} muestre \campointerfaz{Manual}.}

\begin{operacionsensible}
El formato manual define la estructura inicial que usara la competencia. Antes de guardar, revise numero de equipos, cantidad de grupos y clasificados. Para cancelar, cierre la seccion o cambie de pantalla sin presionar \botoninterfaz{Aplicar formato manual}. Para verificar, revise el resumen de la categoria y la distribucion de grupos.
\end{operacionsensible}

\begin{resultadoesperado}
La interfaz confirma que se aplico una configuracion manual para la categoria seleccionada.
\end{resultadoesperado}

\section{Reorganizar grupos}

\figuraManual{capturas/finales/MU-017-distribucion-grupos.png}{Lista de distribucion manual donde cada equipo muestra institucion, grupo actual y selector para moverlo a otro grupo.}{fig:bloque2-mu-017}

\begin{procedimiento}{Ajustar la distribucion manual de grupos}{Operador del torneo}
Use este procedimiento cuando la estructura de grupos existe, pero la organizacion necesita mover equipos antes de operar resultados.
\end{procedimiento}

\begin{precondiciones}
\begin{itemize}
    \item Los equipos ya aparecen en grupos.
    \item La organizacion definio que un movimiento es necesario.
    \item No se han registrado resultados que dependan de la distribucion actual.
\end{itemize}
\end{precondiciones}

\paso{1}{Abra la seccion \campointerfaz{Distribucion manual}.}
\paso{2}{Presione \botoninterfaz{Mostrar lista completa} si la lista esta cerrada.}
\paso{3}{Ubique el robot que desea mover. La fila muestra robot, institucion y grupo actual.}
\paso{4}{En \campointerfaz{Mover a}, seleccione el grupo destino.}
\paso{5}{Repita la revision para cada equipo que deba cambiar de grupo.}
\paso{6}{Presione \botoninterfaz{Guardar distribucion manual}.}
\paso{7}{Espere el mensaje \campointerfaz{Se actualizo la distribucion manual de grupos.} o el mensaje de error correspondiente.}
\paso{8}{Vuelva a abrir la lista y confirme que cada equipo aparece en el grupo esperado.}

\begin{operacionsensible}
Mover equipos cambia la primera fase de competencia. Antes de guardar, revise que ningun grupo quede desbalanceado de forma no autorizada y que no existan resultados previos. Para cancelar, no presione \botoninterfaz{Guardar distribucion manual}. Para verificar, compare la lista despues del guardado con la decision de la organizacion.
\end{operacionsensible}

\begin{fallas}
Si el sistema rechaza la distribucion, revise que todos los destinos pertenezcan a la categoria actual. Si un equipo no aparece, confirme que fue sincronizado al cuadro oficial y que pertenece a la edicion activa.
\end{fallas}

\section{Comprobaciones posteriores}

\begin{procedimiento}{Confirmar que la categoria quedo preparada}{Operador del torneo}
Ejecute esta revision antes de entregar la operacion a la fase de grupos.
\end{procedimiento}

\paso{1}{Revise que el titulo de categoria sea correcto.}
\paso{2}{Confirme el estado de competencia y el modo activo.}
\paso{3}{Verifique que los grupos contienen los equipos esperados.}
\paso{4}{Revise que la cantidad de clasificados por grupo coincide con la decision aprobada.}
\paso{5}{No avance a fases posteriores si hay conteos incompletos, equipos duplicados o dudas de asistencia.}

\begin{resultadoesperado}
La categoria queda lista para operar la fase de grupos en el capitulo correspondiente, sin adelantar resultados ni crear fases posteriores desde esta etapa.
\end{resultadoesperado}

\section{Checklist previo}

\begin{itemize}
    \item Registros revisados en la edicion activa.
    \item Asistencia actualizada o validada por la organizacion.
    \item Equipos completos y sin duplicados evidentes.
    \item Divisiones Colegios y Universidades verificadas por separado.
    \item Permisos de operador confirmados.
    \item Datos de demostracion retirados antes de produccion.
    \item Configuracion inicial revisada antes de aplicar modo sugerido o manual.
    \item Estado previo o progreso existente verificado antes de volver a mezclar.
\end{itemize}

\begin{nota}
La plataforma implementa controles de formato, distribucion y mensajes de resultado. El respaldo operativo previo y la autorizacion para cambiar estructura son recomendaciones de organizacion, no botones automaticos de esta pantalla.
\end{nota}
